\begin{annexitem}{annex:p15.vcl.core}{thm:p15-vcl-core}{Pointwise approximation before taking a supremum}
\begin{proof}
Fix \(C\) and a core sequence with indicators converging pointwise to
\(\mathbf1_C\). Binary convergence means eventual equality at each point.
For this sequence put \(B_k=\bigcup_{j\geq k}(C_j\mathbin{\triangle}C)\).
These are measurable, decrease to the empty set, and have probabilities
tending to zero: the differences \(B_k\setminus B_{k+1}\) are disjoint
and countable additivity expresses \(P(B_k)\) as the tail of a convergent
nonnegative series. Hence
\(|P(C_j)-P(C)|\leq P(B_j)\to0\).
On any fixed finite sample, eventual equality holds simultaneously at
all sampled points, so empirical means also converge. Each full-class
discrepancy is a limit of core discrepancies and is bounded by their
supremum; the opposite inequality follows from inclusion of the core.
The same reasoning on both finite samples proves the two-sample assertion.
Countable suprema of measurable functions are measurable.
Every finite binary pattern is eventually reproduced exactly by the core.
Finally, if \(P_m(C)=0\) and \(P(C)>\varepsilon\), sufficiently late
approximants have the identical empirical pattern and probability still
greater than \(\varepsilon\). Thus a bad consistent member exists iff
one exists in the core, a countable union of measurable events.
\end{proof}
\end{annexitem}

\begin{annexitem}{annex:p15.vcl.sauer}{thm:p15-vcl-sauer}{The trace recursion with all base cases}
\begin{proof}
Prove the bound for any finite family of binary words on \(t\) coordinates
whose shattered coordinate sets have size at most \(d\). An empty family
satisfies the bound trivially. At \(t=0\) there is at most the empty word.
At \(d=0\), two different words would differ at a coordinate and shatter
that singleton, so there is at most one. For \(d\geq t\), use the
trivial bound \(2^t=\sum_{j=0}^t\binom tj\).
In the remaining case project words ending in zero and one to families
\(F_0,F_1\) on \(t-1\) coordinates. Put \(U=F_0\cup F_1\),
\(V=F_0\cap F_1\). Then \(|F|=|U|+|V|\). A coordinate set
shattered by \(U\) is also shattered by \(F\), so \(U\) has dimension
at most \(d\). If \(V\) shatters a set, each of its patterns has
extensions with both last bits. Adjoining the last coordinate therefore
gives a shattered set for \(F\), so \(V\)'s dimension is at most
\(d-1\). Induction on \(t\) and Pascal's identity give
\[
 |F|\leq\sum_{j=0}^d\binom{t-1}j+
             \sum_{j=0}^{d-1}\binom{t-1}j=\sum_{j=0}^d\binom tj,
\]
with out-of-range binomial coefficients zero. For an actual tuple,
shattered coordinate positions must refer to distinct points; its
coordinate dimension is therefore at most the class's VC dimension.
This proves the bound also for tuples with repetitions.

For \(1\leq d\leq t\), set \(x=d/t\). Since \(x^j\geq x^d\)
for \(j\leq d\),
\(x^d\sum_{j=0}^d\binom tj\leq(1+x)^t\leq e^{tx}=e^d\).
The polynomial inequality follows. Dimension zero forces all class
indicators to agree at every point, giving growth one. Infinite dimension
supplies a shattered set at least as large as every prescribed \(t\),
and each subset of a shattered set is shattered, giving growth \(2^t\).
\end{proof}
\end{annexitem}

\begin{annexitem}{annex:p15.vcl.uniform}{thm:p15-vcl-uniform}{Ghost witness, paired random split and conditional finite union}
\begin{proof}
First suppose \(m\varepsilon^2\geq2\). On
\(E=\{\Delta_m>\varepsilon\}\), choose the first core index with
\(|P_m(C)-P(C)|>\varepsilon\). This index is measurable because the
core is countable. For each fixed training sample, its witness is a fixed
event for an independent ghost sample \(Z'_1,\ldots,Z'_m\).
The ghost mean has variance at most \(1/(4m)\); Chebyshev gives
\(P(|P'_m(C)-P(C)|>\varepsilon/2)\leq1/(m\varepsilon^2)\leq1/2\).
With ghost probability at least one half, the two empirical means thus
differ by more than \(\varepsilon/2\). Integrating sections of the product
sample law proves
\[
 P(E)\leq2P\{\sup_C|P_m(C)-P'_m(C)|>\varepsilon/2\}.
\]
No unmeasurable maximizing selection was used; a supremum strictly above
the threshold has a countable-core witness even when it is not attained.

Fix the pooled pairs \((z_i,z'_i)\), and draw independent fair signs
\(\sigma_i\). A positive sign retains the pair's order and a negative
sign swaps it. Because the original pairs are iid with identical laws
in both coordinates, the unconditional law of this randomized split is
again the original two-sample law. For a fixed trace its signed difference
is \(m^{-1}\sum_i\sigma_i a_i\), with
\(a_i=\mathbf1_C(z_i)-\mathbf1_C(z'_i)\in\{-1,0,1\}\).
For real \(u\), the power series and \((2k)!\geq2^kk!\) give
\(\cosh u\leq e^{u^2/2}\). Independence therefore gives
\(E_\sigma e^{\lambda\sum_i\sigma_i a_i}
 \leq e^{\lambda^2\sum_i a_i^2/2}\).
Markov's inequality and optimization at
\(\lambda=(m\varepsilon/2)/\sum_i a_i^2\) bound one tail by
\(e^{-m\varepsilon^2/8}\), since \(\sum_i a_i^2\leq m\).
The negative tail has the same bound. If all coefficients vanish,
the tail event is empty and no division is performed.

Only \(\Pi_{\mathcal C}(2m)\) distinct pooled traces can occur.
Conditionally, their event probabilities are summed to give at most
\(2\Pi_{\mathcal C}(2m)e^{-m\varepsilon^2/8}\).
This conditional event is measurable: it is also a union over the fixed
countable core, and, for each pool, identical traces give identical sign
events. One may choose the first core representative of each trace; its
index conditions are countable measurable comparisons. Integrating and
including the ghost factor two proves the factor four.
For \(m\varepsilon^2<2\), the expression before clipping exceeds
\(4e^{-1/4}>1\), so the trivial bound one gives the asserted minimum.
For dimension zero use the one fixed indicator and the already proved
Hoeffding bound of Theorem~\ref{thm:p15-conc-exponential}.
For fixed finite \(d\), Sauer gives polynomial growth; its logarithm
is \(O_d(\ln m)\), dominated by \(m\varepsilon^2/8\).
The constants are independent of \(P\), proving the stated uniform limit.
\end{proof}
\end{annexitem}

\begin{annexitem}{annex:p15.vcl.erm}{thm:p15-vcl-erm}{Loss traces, measurable minimization and holdout}
\begin{proof}
On a fixed labelled tuple, the error vector is the prediction vector
xor its fixed label vector. This is a bijection on binary vectors, so
the number of traces equals the number of prediction traces. Maximizing
over labelled tuples gives no larger growth than \(\Pi_H\); taking
all labels zero gives the reverse inequality. Pointwise core approximation
of predictions also approximates their binary losses on every labelled
point, so the preceding measurability result applies to the loss class.

Empirical errors take values in the finite set
\(\{0,1/m,\ldots,1\}\). The nonempty set of attained values has
a minimum, and every full-class trace occurs in the core. Thus some core
index attains that same minimum. The infimum of the countable measurable
core scores is measurable. The event that index \(j\) is its first
attainer is equality to that infimum intersected with failure of all
earlier indices to attain it. These measurable events partition sample
space, defining a measurable index. Evaluation of its selected predictor
is jointly measurable by the countable partition. This constructs the
selection but does not provide a finite algorithm for detecting that no
earlier unexamined hypothesis has a smaller score.

On uniform discrepancy at most \(u\), for every comparator \(h\),
\[
 R(\widehat h)\leq\widehat R_m(\widehat h)+u
 \leq\widehat R_m(h)+\eta+u\leq R(h)+2u+\eta.
\]
Taking the infimum proves the excess-risk bound without a population
minimizer assumption. In a realizable problem, a risk-zero member makes
no sample errors almost surely, by a finite union of its null error events.
The empirical minimum is consequently zero, including when that member
is represented in the core only through pointwise approximation.

For holdout, write the joint law as the training law times \(P^{\otimes v}\).
For every fixed training object its selected loss is one measurable
Bernoulli loss on independent validation records. Hoeffding gives
\(2e^{-2vu^2}\). Joint measurability permits integration over training,
giving the same bound. This argument requires no regular conditional
distribution and does not permit selecting the predictor using the holdout.
\end{proof}
\end{annexitem}

\begin{annexitem}{annex:p15.vcl.consistent}{thm:p15-vcl-consistent}{The consistent ghost count and random-split probability}
\begin{proof}
The bad-consistent event is a countable-core union by the core lemma.
Choose its first witness and write its true error as \(p>\varepsilon\).
On the independent ghost sample its number of errors \(B\) has mean
\(mp\) and variance at most \(mp\). Chebyshev bounds
\(P(B<mp/2)\leq4/(mp)<4/(m\varepsilon)\leq1/2\).
Thus the bad event's probability is at most twice that of a member
having zero training errors and at least \(m\varepsilon/2\) ghost errors.

Use the same fixed-pool paired random split. For a fixed loss trace,
a pair with two errors makes zero training error impossible. Otherwise
let \(k\) pairs have exactly one error. To avoid every training error,
each of those pairs must orient its one error to the ghost side, with
probability \(2^{-k}\); all other pairs have zero errors. The ghost
then has exactly \(k\) errors. Requiring \(k\geq m\varepsilon/2\)
gives probability at most \(2^{-m\varepsilon/2}\), with no integer
rounding loss. Sum over at most \(\Pi_H(2m)\) loss traces and include
the ghost factor two. The same core argument makes this conditional
finite union measurable. In a realizable problem a consistent output
exists almost surely by the ERM construction; under noise that premise
is not automatic. In dimension zero all predictors agree pointwise,
and a risk-zero member therefore makes every member risk zero.
\end{proof}
\end{annexitem}

\begin{annexitem}{annex:p15.vcl.samples}{thm:p15-vcl-samples}{Explicit sufficient bounds and the ceiling step}
\begin{proof}
For the uniform bound put \(A=\ln(64/\varepsilon^2)\),
\(B=\ln(8/\delta)\), \(L=dA+B\) and \(m_0=64L/\varepsilon^2\).
Here \(A>4\), \(m_0\geq d\) and \(m_0\varepsilon^2\geq2\).
Sauer reduces the log of the unclipped failure bound to
\(\ln4+d\ln(2em/d)-m\varepsilon^2/8\).
At \(m_0\), the required upper exponent terms satisfy
\[
 d\ln(2em_0/d)+\ln(4/\delta)
 =d[A+\ln(2e)+\ln(L/d)]+\ln(4/\delta)
 \leq2L+2d<3L<8L=m_0\varepsilon^2/8.
\]
We used \(\ln x\leq x\), \(\ln(2e)<2\), and
\(\ln(4/\delta)<B\). The difference
\(m\varepsilon^2/8-d\ln(2em/d)\) is increasing for
\(m\geq8d/\varepsilon^2\), already true at \(m_0\).
Therefore taking a ceiling, or any larger integer, preserves the bound.

For consistent learning put \(A=\ln(64/\varepsilon)\),
\(B=\ln(4/\delta)\), \(L=dA+B\), \(m_0=32L/\varepsilon\).
It satisfies \(m_0\varepsilon\geq8\) and \(m_0\geq d\).
Since \(\ln2\geq1/2\), it suffices to control
\(2(2em/d)^d e^{-m\varepsilon/4}\). At \(m_0\),
\[
 d\ln(2em_0/d)+\ln(2/\delta)
 =d[A+1+\ln(L/d)]+\ln(2/\delta)
 \leq2L+d<3L<8L=m_0\varepsilon/4.
\]
The exponent difference is increasing for \(m\geq4d/\varepsilon\),
so the ceiling and all larger integers also suffice. Finally exact ERM
has excess at most twice the uniform discrepancy, justifying the
substitution \(\varepsilon=\tau/2\). These computations are sufficient
upper bounds, not optimality claims.
\end{proof}
\end{annexitem}

\begin{annexitem}{annex:p15.vcl.examples}{thm:p15-vcl-examples}{Thresholds, capacity failure, memorization and irreducible noise}
\begin{proof}
For any real threshold choose rational thresholds decreasing to it.
Their closed-threshold indicators converge pointwise, including at the
threshold itself. Constants belong to the core as well. On \(t\) distinct
ordered points only the \(t+1\) initial-segment patterns occur, and every
one is attained. A singleton is shattered but two distinct points cannot
have the pattern zero at the smaller and one at the larger. The dimension
is therefore one; repetitions cannot add patterns.

Finite subsets of \(\mathbb N\) are countable, by encoding their finite
increasing lists, and they realize every pattern on every finite subset.
For the uniform law on \(2m\) points, let \(A\) be the finite set of
observed points. Then \(P_m(A)=1\), whereas \(P(A)\leq m/(2m)=1/2\).
Thus discrepancy is at least one half pathwise. Under a point mass every
sample point equals its atom almost surely and every event's empirical
frequency equals its probability, giving zero discrepancy. The different
laws show why distribution-free statements cannot be read as necessity
for each fixed law.

For the uniform-real memorizer, duplicate training points have the same
deterministic target label. The classifier has that label at each training
point and predicts one elsewhere, so its empirical loss is zero. Its
evaluation is jointly Borel measurable: matching a given real point to
one of finitely many sample coordinates is a finite union of measurable
diagonal events, with the stated finite label rule. Its negative predictions
are only at finitely many sampled points below one half. That set has
uniform probability zero. It therefore errs on almost every point below
one half and nowhere above, giving risk one half. This calculation concerns
a specific measurable selected classifier, not a presumption that every
uncountable-class supremum is measurable.
Finally, if the label is an independent fair bit, for any measurable
predictor
\(R(h)=\tfrac12P(h(X)=1)+\tfrac12P(h(X)=0)=1/2\).
Thus no member is risk zero, while the minimum risk within any nonempty
deterministic class is still one half. Excess-risk control and realizability
are different claims.
\end{proof}
\end{annexitem}
